%% notes-tex/030-scaling-plots/sections/5-pair.tex

\section{The 025 pair\secstatus{todo}}
\label{sec:pair}

The 025 twin trains the merged-label trainer on the same subsets, the same
genotypes pair for pair, at 180x8, with three seeds at the two smallest
subsets. It passes through the same readout and the same two reads at the same
$B$, so the two plots of Figure~\ref{fig:pair} compare like with like.

\tcfig{figures/030-scaling-plots-pair.pdf}{%
The pair family, sketched. \textbf{a}, P1: the 030 minus 025 pooled validation
MSE per subset, seed mean and sd where both sides have three seeds, a single
point at 1/4, and the full point hollow as the difference of two unconverged
ten-epoch runs. \textbf{b}, P2: the paired learning curves against records seen,
one panel per subset in the real figure; a gap that holds past $B$ is more
supervision per genotype, one that closes is a faster start.
}{fig:pair}
\src{notes/assets/drawio/030-scaling-plots.gen.py}

\notesec{P1, the gap against data}
Pooled only, since 025 logs no per-screen key. Interpretable at 1/64 and 1/16,
where both sides converge and both have a seed spread; a single comparison at
1/4; a learning-speed difference at the full pool. The gap is per-entry rows
plus the dataset token plus the measurement types the 030 build added, and the
caption says so until the token ablation column exists.

\notesec{P2, paired learning curves}
The plot that tells the two readings of a gap apart. If the 030 curve sits below
the 025 curve from the first epochs and the distance holds through $B$ and
beyond, the per-entry rows are more supervision at every stage of training. If
the curves start apart and meet, the extra rows were a faster start that the
merged label catches up on, and the gap at a fixed $B$ overstates what the rows
buy at convergence.
